% Figure: hydrostatic force on a curved (quarter-circle) gate, decomposed into a
% horizontal component (force on the vertical projection) and a vertical
% component (weight of the fluid volume above the surface). No em-dashes.
\begin{figure}[htbp]
\centering
\begin{tikzpicture}[scale=1.1]
% free surface
\draw[engblue, thick] (-0.6,3.2) -- (4.4,3.2);
\node[engblue, font=\scriptsize, anchor=south] at (-0.4,3.2) {free surface};
% water region
\fill[engblue!12] (-0.6,0) rectangle (4.4,3.2);
% quarter-circle curved gate: centre at (3.0,3.2), radius R=2.4, sweeping the
% lower-left quadrant from (0.6,3.2) down to (3.0,0.8)
\def\cx{3.0}\def\cy{3.2}\def\R{2.4}
\draw[engblack, ultra thick] (\cx,\cy-\R) arc (270:180:\R);
\node[engblack, font=\small, anchor=west] at (1.0,1.4) {curved gate};
% the fluid volume directly above the curved surface (for the vertical component)
\fill[englime!25] (\cx-\R,\cy) -- (\cx-\R,\cy) arc (180:270:\R) -- (\cx,\cy) -- cycle;
% vertical projection of the curved surface (a vertical line of height R)
\draw[engamber, thick, dashed] (\cx-\R,\cy) -- (\cx-\R,\cy-\R);
% horizontal component arrow (acts on the vertical projection)
\draw[engamber, ultra thick, ->] (-0.2,1.6) -- (0.6,1.6);
\node[engamber, font=\small, anchor=east] at (-0.2,1.6) {$F_H$};
% vertical component arrow (weight of fluid above, acting up on the gate)
\draw[englime, ultra thick, ->] (2.1,1.0) -- (2.1,0.1);
\node[englime, font=\small, anchor=west] at (2.15,0.55) {$F_V$};
% resultant
\draw[engred, ultra thick, ->] (1.35,1.3) -- ($(1.35,1.3)+(0.75,-0.78)$);
\node[engred, font=\small, anchor=north east] at ($(1.35,1.3)+(0.75,-0.78)$) {$F=\sqrt{F_H^2+F_V^2}$};
\end{tikzpicture}
\caption[Curved surface by the component method]{Hydrostatic force on a
curved gate, resolved by the horizontal-and-vertical component method. The
horizontal component $F_H$ (amber) equals the hydrostatic force on the
flat vertical projection of the curved surface, computed by the ordinary
$p_CA$ rule. The vertical component $F_V$ (green) equals the weight of the
fluid volume lying between the curved surface and the free surface (shaded);
here the water sits above the gate, so $F_V$ presses down, and by reaction
the gate pushes up on the water. The two components combine vectorially
into the resultant $F$ (red), whose line of action passes through the
centre of the arc because every pressure on a circular surface points along
a radius.}
\label{fig:vol03ch10:curved-surface}
\end{figure}
